\documentclass[12pt, twoside]{article}
%\documentclass[12pt, twoside]{article}
\usepackage[letterpaper, margin=1in, headsep=0.2in]{geometry}
\setlength{\headheight}{0.6in}
%\usepackage[english]{babel}
\usepackage[utf8]{inputenc}
\usepackage{microtype}
\usepackage{amsmath}
\usepackage{amssymb}
%\usepackage{amsfonts}
\usepackage[nomessages]{fp} %\FPeval{\var-name}{2*sin(pi/6)}
\usepackage{siunitx} %units in math. eg 20\milli\meter
\usepackage{yhmath} % for arcs, overparenth command
\usepackage{tikz} %graphics
\usetikzlibrary{quotes, angles, arrows, arrows.meta}
\usepackage{pgfplots}
\usepgfplotslibrary{fillbetween}
\pgfplotsset{compat=1.16}
\usepackage{graphicx} %consider setting \graphicspath{{images/}}
\usepackage{parskip} %no paragraph indent
\usepackage{enumitem}
\usepackage{multicol}
\usepackage{venndiagram}

\usepackage{fancyhdr}
\pagestyle{fancy}
\fancyhf{}
\renewcommand{\headrulewidth}{0pt} % disable the underline of the header
\raggedbottom
\hfuzz=2mm %suppresses overfull box warnings

\usepackage{hyperref}
\usepackage{eurosym}

\fancyhead[LE]{\thepage}
\fancyhead[RO]{Markscheme}
\fancyhead[LO]{La Scuola d'Italia / Dr.\ Huson / IB Math 11 \\* 13 October 2026}

\begin{document}
\subsubsection*{1.11 Classwork: Geometric Series --- Markscheme}

\begin{enumerate}[itemsep=0.3cm]

%--- Problem 1: S_5 by adding, S_5 by formula, S_10, n from S_n [7 marks] ---
\item[1.] [Maximum mark: 7]

The series $2 + 6 + 18 + 54 + \cdots$, so $u_{1} = 2$ and $r = 3$.

\medskip

\textbf{(a)} $S_{5} = 2 + 6 + 18 + 54 + 162 = 242$ \hfill \textbf{A1}

\emph{Accept $242$ obtained by any method, including the sum formula.}

\emph{The sum of only four terms, $2 + 6 + 18 + 54 = 80$, earns no mark.
A fifth term from adding (e.g.\ $54 + 36 = 90$) in place of multiplying by
$3$ earns no mark.}

\medskip

\textbf{(b)} $S_{5} = \dfrac{2\left(3^{5} - 1\right)}{3 - 1}$ \hfill \textbf{M1}

$S_{5} = \dfrac{2(242)}{2} = 242$ \hfill \textbf{A1}

\emph{Accept substitution into the second formula,
$S_{5} = \dfrac{2\left(1 - 3^{5}\right)}{1 - 3} = \dfrac{-484}{-2} = 242$.}

\emph{M1 is for substituting $u_{1} = 2$, $r = 3$ and $n = 5$ into a sum
formula, however the arithmetic is then arranged. The question asks for
the substituted values, so $242$ with no substitution shown earns the A1
only.}

\emph{An exponent of $4$ in place of $5$, giving $80$, earns the M1 only.}

\medskip

\textbf{(c)} $S_{10} = \dfrac{2\left(3^{10} - 1\right)}{3 - 1}$ \hfill \textbf{M1}

$S_{10} = 59\,048$ \hfill \textbf{A1}

\emph{An exponent of $9$, giving $19\,682$, earns the M1 only.
$u_{10} = 2 \times 3^{9} = 39\,366$ (a term, not a sum) earns no marks.}

\emph{Accept $59\,048$ from a GDC sum of the first $10$ terms for full
marks. Feedback: write out the formula substitution so the method is
visible.}

\medskip

\textbf{(d)} $\dfrac{2\left(3^{n} - 1\right)}{3 - 1} = 177\,146$ \hfill \textbf{M1}

$3^{n} - 1 = 177\,146 \ \Rightarrow \ 3^{n} = 177\,147$

$n = 11$ \hfill \textbf{A1}

\emph{Accept solving the equation on the GDC, with logarithms, or by a
table of sums ($S_{10} = 59\,048$, $S_{11} = 177\,146$). Accept trial of
powers of $3$ ($3^{11} = 177\,147$).}

\emph{M1 is for an equation setting a sum expression in $n$ equal to
$177\,146$, or for a table of sums reaching $177\,146$. Setting a term
expression equal to $177\,146$ ($2 \times 3^{\,n-1} = 177\,146$) earns no
marks.}

\newpage

%--- Problem 2: Starting salaries rising 6% per year (NYC $70 000, Italy 30 000 euros); 3rd-year salary, 3-year total, first total over 1 000 000 for each [8 marks] ---
\item[2.] [Maximum mark: 8]

Luca: $u_{1} = 70\,000$ dollars; Sofia: $u_{1} = 30\,000$ euros; both
$r = 1.06$.

\medskip

\textbf{(a)} $u_{3} = 70\,000(1.06)^{2}$ \hfill \textbf{M1}

$=$ \$$78\,652$ \hfill \textbf{A1}

\emph{Accept $78\,700$ (3 s.f.). Accept repeated multiplication:
$70\,000 \times 1.06 = 74\,200$, then $74\,200 \times 1.06 = 78\,652$.}

\emph{An exponent of $3$ (giving \$$83\,371.12$) earns the M1 only. A
multiplier of $0.06$ or $6$ earns no marks.}

\medskip

\textbf{(b)} $S_{3} = \dfrac{70\,000\left(1.06^{3} - 1\right)}{1.06 - 1}$ \hfill \textbf{M1}

$=$ \$$222\,852$ \hfill \textbf{A1}

\emph{Accept adding the three salaries,
$70\,000 + 74\,200 + 78\,652 = 222\,852$, for full marks. Accept $223\,000$
(3 s.f.).}

\emph{FT: award M1A1 for the sum of the first three salaries using the
candidate's own third-year salary from (a). Three times the first salary
(\$$210\,000$) or the third-year salary alone earns no marks.}

\medskip

\textbf{(c)} $\dfrac{70\,000\left(1.06^{n} - 1\right)}{1.06 - 1} > 1\,000\,000$ \hfill \textbf{M1}

After $10$ years: \$$922\,655.65$ \ ($< 1\,000\,000$); \quad after $11$
years: \$$1\,048\,014.98$ \ ($> 1\,000\,000$)

$11$ years \hfill \textbf{A1}

\emph{Accept solving the inequality (or the equation) on the GDC or with
logarithms: $1.06^{n} > \frac{13}{7}$, $n > 10.6\ldots$, so $11$ years.}

\emph{Accept a table search: full marks for showing both adjacent rows
($n = 10$ and $n = 11$) and stating $11$. A single row still earns full
marks, with the feedback note that the adjacent row $n = 10$ is what
justifies ``first''. See markscheme-spec §18.}

\emph{$10.6$ is not an acceptable final answer --- the answer must be a
whole number of years. Rounding down to $10$ earns the M1 only. Using the
yearly salary $70\,000(1.06)^{\,n-1}$ in place of the total earns no marks.}

\medskip

\textbf{(d)} $\dfrac{30\,000\left(1.06^{n} - 1\right)}{1.06 - 1} > 1\,000\,000$ \hfill \textbf{M1}

After $18$ years: \euro$927\,169.58$ \ ($< 1\,000\,000$); \quad after $19$
years: \euro$1\,012\,799.75$ \ ($> 1\,000\,000$)

$19$ years \hfill \textbf{A1}

\emph{Accept the GDC, logarithms ($1.06^{n} > 3$, $n > 18.9$, so $19$ years)
or a table search, on the same terms as (c). Rounding down to $18$ earns the
M1 only.}

\end{enumerate}
\end{document}
